\documentclass[12pt,letterpaper, onecolumn]{exam}
\usepackage[utf8]{inputenc}
\usepackage[T1]{fontenc}
\usepackage{amsmath}
\usepackage{amssymb}
\usepackage[lmargin=71pt, tmargin=1.2in]{geometry}
\usepackage{xcolor}
\usepackage[brazilian]{babel}
\usepackage{natbib}
\usepackage{enumitem}
\usepackage{booktabs}
\usepackage[hidelinks]{hyperref}
\urlstyle{same}

\renewenvironment{solution}
{\par\noindent\textbf{R:}\enspace\color{blue}}
{\par\color{black}}

\thispagestyle{empty}

\begin{document}

\begingroup
\centering
\LARGE Lista 9 Gabarito - Econometria I - 2026\\[0.5em]
\large Prof: André Luiz Pereira Mancha \par
\large Monitor: Tuffy Licciardi Issa \par
\endgroup
\rule{\textwidth}{0.4pt}
\pointsdroppedatright
\printanswers
\renewcommand{\solutiontitle}{\noindent\textbf{Ans:}\enspace}
\newcommand{\qstar}{\hspace{-0.6em}\textsuperscript{*}\hspace{0.2em}}

\begin{questions}

\question (7.10, \citep{wooldridge2016}) Para uma criança \(i\) que mora em determinada região de ensino, defina \textit{voucher} como uma variável \textit{dummy} igual a um se a criança foi selecionada para participar de um programa de bolsas de estudos em uma escola, e defina \textit{score} como a nota da criança em um exame padronizado subsequente. Suponha que a variável de participação, \textit{voucher}, seja completamente aleatorizada para que ela seja independente tanto dos fatores observados quanto dos não observados que possam afetar a nota do teste de avaliação.
\begin{enumerate}[label=(\roman*)]
    \item Se você executar uma regressão simples de \textit{score} sobre \textit{voucher}, usando uma amostra aleatória de tamanho \(n\), o estimador de MQO produziria um estimador não viesado do efeito do programa de bolsas de estudos?
    \item Suponha que você possa coletar informações adicionais sobre perfis familiares, tais como renda familiar, estrutura familiar (por exemplo, se a criança mora com os dois pais) e nível de escolaridade dos pais. Você precisaria controlar esses fatores para obter um estimador não viesado dos efeitos do programa de bolsas de estudos? Explique.
    \item Por que você precisaria incluir as variáveis de perfis familiares na regressão? Existe alguma situação em que você não incluiria essas variáveis?
\end{enumerate}

\begin{solution}
\begin{enumerate}[label=(\roman*)]
    \item Sim. Como \textit{voucher} foi completamente aleatorizada, ela é independente dos fatores observados e não observados que afetam \textit{score}. Assim, na regressão simples, a condição de exogeneidade é satisfeita e o estimador de MQO é não viesado para o efeito causal médio do programa.

    \item Não para obter não viesamento. A aleatorização já garante identificação causal sem esses controles. Ainda assim, incluir variáveis familiares pode reduzir a variância do erro e melhorar a precisão das estimativas.

    \item O principal motivo para incluí-las é ganho de precisão. Se renda, estrutura familiar e escolaridade dos pais explicam parte importante de \textit{score}, então controlá-las reduz a variância residual. Uma situação em que eu poderia não incluí-las é quando a aleatorização é confiável e essas covariáveis pouco acrescentam em termos de precisão, ou quando coletá-las é custoso e pouco informativo.
\end{enumerate}
\end{solution}

\vspace{0.7cm}

\question (7.9, \citep{wooldridge2016}) Seja \(d\) uma variável \textit{dummy} (binária) e \(z\) uma variável quantitativa. Considere o modelo
\[
y = \beta_0 + \delta_0 d + \beta_1 z + \delta_1 d \cdot z + u;
\]
essa é uma versão geral de um modelo com uma interação entre uma variável \textit{dummy} e uma variável quantitativa. [Um exemplo está na equação (7.17).]
\begin{enumerate}[label=(\roman*)]
    \item Como isso não alterará nada importante, defina o erro com valor zero, \(u=0\). Então, quando \(d=0\), podemos escrever o relacionamento entre \(y\) e \(z\) como a função \(f_0(z)=\beta_0+\beta_1 z\). Escreva essa mesma relação quando \(d=1\), em que você deve usar \(f_1(z)\) no lado esquerdo para denotar a função linear de \(z\).
    \item Considerando \(\delta_1 \neq 0\) (o que significa que as duas não são paralelas), demonstre que o valor de \(z^\ast\) é tal que \(f_0(z^\ast)=f_1(z^\ast)\) e \(z^\ast = -\delta_0/\delta_1\). Esse é o ponto no qual as duas linhas se cruzam [como na Figura 7.2(b)]. Demonstre que \(z^\ast\) será positivo se, e somente se, \(\delta_0\) e \(\delta_1\) tiverem sinais opostos.
    \item Usando os dados contidos no arquivo \texttt{TWOYEAR}, a seguinte equação poderá ser estimada:
    \[
    \widehat{\log(wage)} = 2{,}289 - 0{,}357\,female + 0{,}050\,totcoll + 0{,}030\,female \cdot totcoll
    \]
    \[
    \hspace{2.0cm} (0{,}011)\hspace{1.2cm}(0{,}015)\hspace{1.1cm}(0{,}003)\hspace{1.0cm}(0{,}005)
    \]
    \[
    n = 6{,}763,\; R^2 = 0{,}202
    \]
    em que todos os coeficientes e erros padrão foram arredondados com três casas decimais. Usando essa equação, encontre o valor de \textit{totcoll} de tal forma que os valores previstos de \(\log(wage)\) sejam os mesmos tanto para homens quanto para mulheres.
    \item Com base na equação do item (iii), as mulheres podem, de forma realística, obter educação superior de modo que seus ganhos alcancem os dos homens? Explique.
\end{enumerate}

\begin{solution}
\begin{enumerate}[label=(\roman*)]
    \item Quando \(d=1\),
    \[
    y = \beta_0 + \delta_0 + (\beta_1+\delta_1)z,
    \]
    de modo que
    \[
    f_1(z)= (\beta_0+\delta_0)+(\beta_1+\delta_1)z.
    \]

    \item O ponto de cruzamento satisfaz
    \[
    f_0(z^\ast)=f_1(z^\ast).
    \]
    Logo,
    \[
    \beta_0+\beta_1 z^\ast = \beta_0+\delta_0+(\beta_1+\delta_1)z^\ast.
    \]
    Cancelando os termos comuns,
    \[
    0=\delta_0+\delta_1 z^\ast,
    \]
    e portanto
    \[
    z^\ast = -\frac{\delta_0}{\delta_1}.
    \]
    Além disso, \(z^\ast>0\) se, e somente se, \(\delta_0\) e \(\delta_1\) tiverem sinais opostos.

    \item Para homens (\(female=0\)),
    \[
    \widehat{\log(wage)} = 2{,}289 + 0{,}050\,totcoll.
    \]
    Para mulheres (\(female=1\)),
    \[
    \widehat{\log(wage)} = 2{,}289 - 0{,}357 + (0{,}050+0{,}030)\,totcoll
    = 1{,}932 + 0{,}080\,totcoll.
    \]
    Igualando as duas expressões:
    \[
    2{,}289 + 0{,}050\,totcoll = 1{,}932 + 0{,}080\,totcoll.
    \]
    Assim,
    \[
    0{,}357 = 0{,}030\,totcoll
    \quad \Rightarrow \quad
    totcoll = 11{,}9.
    \]

    \item Em termos práticos, não parece muito realístico. O valor de \(11{,}9\) anos de ensino superior é muito alto: corresponde a uma trajetória muito longa após o ensino médio. Portanto, embora o modelo permita esse cruzamento, ele parece ocorrer em uma faixa pouco plausível da amostra.
\end{enumerate}
\end{solution}

\vspace{0.7cm}

\question (7.11, \citep{wooldridge2016}) As equações a seguir foram estimadas usando os dados do arquivo \texttt{ECONMATH}, com erros padrão registrados abaixo dos coeficientes. A nota média da classe, medida como porcentagem, é de cerca de \(72{,}2\); exatamente \(50\%\) dos estudantes são do sexo masculino; e a média de \textit{colgpa} (nota média no início do semestre) é de cerca de \(2{,}81\).

\[
\widehat{score} = 32{,}31 + 14{,}32\,colgpa
\]
\[
\hspace{2.2cm} (2{,}00)\hspace{1.4cm}(0{,}70)
\]
\[
n = 856,\; R^2 = 0{,}329,\; \bar{R}^2 = 0{,}328
\]

\[
\widehat{score} = 29{,}661 + 3{,}83\,male + 14{,}57\,colgpa
\]
\[
\hspace{2.2cm} (2{,}04)\hspace{1.4cm}(0{,}74)\hspace{1.4cm}(0{,}69)
\]
\[
n = 856,\; R^2 = 0{,}349,\; \bar{R}^2 = 0{,}348
\]

\[
\widehat{score} = 30{,}36 + 2{,}47\,male + 14{,}33\,colgpa + 0{,}479\,male \cdot colgpa
\]
\[
\hspace{2.2cm} (2{,}86)\hspace{1.2cm}(3{,}96)\hspace{1.2cm}(0{,}98)\hspace{1.2cm}(1{,}383)
\]
\[
n = 856,\; R^2 = 0{,}349,\; \bar{R}^2 = 0{,}347
\]

\[
\widehat{score} = 30{,}36 + 3{,}82\,male + 14{,}33\,colgpa + 0{,}479\,male \cdot (colgpa - 2{,}81)
\]
\[
\hspace{2.2cm} (2{,}86)\hspace{1.2cm}(0{,}74)\hspace{1.2cm}(0{,}98)\hspace{1.2cm}(1{,}383)
\]
\[
n = 856,\; R^2 = 0{,}349,\; \bar{R}^2 = 0{,}347
\]

\begin{enumerate}[label=(\roman*)]
    \item Interprete o coeficiente sobre \textit{male} na segunda equação e construa um intervalo de confiança de \(95\%\) para \(\beta_{male}\). Esse intervalo de confiança exclui zero?
    \item Na segunda equação, por que a estimativa de \textit{male} é tão imprecisa? Devemos agora concluir que não existem diferenças de gênero em notas depois de controlar \textit{colgpa}? [Dica: Você pode querer calcular a estatística \(F\) da hipótese nula de que não há diferença de gênero no modelo com interação.]
\end{enumerate}

\begin{solution}
\begin{enumerate}[label=(\roman*)]
    \item Na segunda equação, o coeficiente de \textit{male} é \(3{,}83\). Portanto, mantendo \textit{colgpa} fixa, homens obtêm em média \(3{,}83\) pontos a mais do que mulheres. O intervalo de confiança de \(95\%\) é
    \[
    3{,}83 \pm 1{,}96(0{,}74),
    \]
    isto é,
    \[
    [2{,}38,\; 5{,}28].
    \]
    Como zero não pertence a esse intervalo, ele exclui zero.

    \item A dica do enunciado sugere que a preocupação real está no modelo com interação. Nele, o coeficiente de \textit{male} sem centragem mede a diferença de gênero quando \textit{colgpa}\(=0\), um valor muito distante da média observada (\(2{,}81\)). Por isso o erro padrão fica grande nessa parametrização. Não devemos concluir automaticamente que não há diferenças de gênero apenas porque um coeficiente isolado parece impreciso. O teste correto, no modelo com interação, é a hipótese conjunta
    \[
    H_0:\beta_{male}=0 \quad \text{e} \quad \beta_{male\cdot colgpa}=0.
    \]
    Portanto, a conclusão deve se basear em um teste \(F\) conjunto, e não apenas no \(t\) de um coeficiente.
\end{enumerate}
\end{solution}

\vspace{0.7cm}

\question (7.12, \citep{wooldridge2016}) Consideremos o Exemplo 7.11, em que, antes de calcular a interação entre a raça/etnicidade de um jogador e a composição racial da cidade, centramos as variáveis de composição da cidade sobre as médias amostrais, \textit{percblck} e \textit{perchisp} (que são, aproximadamente, \(16{,}55\) e \(10{,}82\), respectivamente). A equação estimada resultante é

\[
\widehat{\log(salary)} = 10{,}23 + 0{,}0673\,years + 0{,}0089\,gamesyr + 0{,}00095\,bavg + 0{,}146\,hrunsyr
\]
\[
\hspace{1.8cm} (2{,}18)\hspace{1.0cm}(0{,}0129)\hspace{0.9cm}(0{,}0034)\hspace{0.9cm}(0{,}00151)\hspace{0.9cm}(0{,}164)
\]
\[
\hspace{0.3cm} +\; 0{,}0045\,rbisyr + 0{,}0072\,runsyr + 0{,}0011\,fldperc + 0{,}0075\,allstar
\]
\[
\hspace{1.8cm} (0{,}0076)\hspace{0.9cm}(0{,}0046)\hspace{0.9cm}(0{,}0021)\hspace{0.9cm}(0{,}0029)
\]
\[
\hspace{0.7cm} +\; 0{,}0080\,black + 0{,}0273\,hispan + 0{,}0125\,black \cdot (\textit{percblck} - \overline{\textit{percblck}})
\]
\[
\hspace{1.8cm} (0{,}0840)\hspace{1.0cm}(0{,}1084)\hspace{1.0cm}(0{,}0050)
\]
\[
\hspace{0.7cm} +\; 0{,}0201\,hispan \cdot (\textit{perchisp} - \overline{\textit{perchisp}})
\]
\[
\hspace{1.8cm} (0{,}0098)
\]
\[
n = 330,\; R^2 = 0{,}638
\]

\begin{enumerate}[label=(\roman*)]
    \item Por que os coeficientes de \textit{black} e \textit{hispan} são agora tão diferentes daqueles que foram reportados na equação (7.19)? Especificamente, como se pode interpretar esses coeficientes?
    \item O que você pensa do fato de nem \textit{black} nem \textit{hispan} serem estatisticamente significantes na equação acima?
    \item Ao comparar a equação acima com a equação (7.19), mais alguma coisa mudou? Por que sim ou por que não?
\end{enumerate}

\begin{solution}
\begin{enumerate}[label=(\roman*)]
    \item Eles mudam porque, ao centrar \textit{percblck} e \textit{perchisp}, os coeficientes de \textit{black} e \textit{hispan} passam a medir o efeito de raça/etnia quando a composição racial da cidade está em seu valor médio amostral. Portanto, \(0{,}0080\) é o diferencial salarial estimado para jogadores negros em uma cidade com \textit{percblck} média, e \(0{,}0273\) é o diferencial para jogadores hispânicos em uma cidade com \textit{perchisp} média.

    \item Isso não é especialmente preocupante. Como os efeitos de raça/etnia dependem da composição racial da cidade, o efeito avaliado exatamente na média pode ser pequeno e impreciso, mesmo quando os termos de interação são relevantes. O ponto central é que o diferencial não é constante: ele varia com o contexto da cidade.

    \item Sim e não. Os coeficientes de \textit{black} e \textit{hispan} mudam de interpretação, e o intercepto também muda. Mas os valores ajustados, os resíduos, o \(R^2\), o ajuste geral e os coeficientes das interações permanecem essencialmente os mesmos. Em outras palavras, houve uma reparametrização, não uma mudança substantiva do modelo.
\end{enumerate}
\end{solution}

\vspace{0.7cm}

\question\qstar (7.C14, \citep{wooldridge2016}) Use os dados do arquivo \texttt{CHARITY} para responder a essa questão. A variável \textit{respond} é uma variável \textit{dummy} igual a um se a pessoa respondeu com uma contribuição à mais recente correspondência enviada por uma organização de caridade. A variável \textit{resplast} é uma variável \textit{dummy} igual a um se a pessoa respondeu ao envio anterior, \textit{avggift} é a média de doações passadas (em florins holandeses), e \textit{propresp} é a proporção de vezes que a pessoa respondeu às correspondências anteriores.
\begin{enumerate}[label=(\roman*)]
    \item Estime um modelo de probabilidade linear relacionando \textit{respond} a \textit{resplast} e \textit{avggift}. Registre os resultados na forma usual e interprete o coeficiente sobre \textit{resplast}.
    \item O valor médio das doações passadas parece afetar a probabilidade de resposta?
    \item Adicione a variável \textit{propresp} ao modelo e interprete seu coeficiente. [Tenha cuidado aqui: o aumento de uma unidade em \textit{propresp} é a maior mudança possível.]
    \item O que aconteceu com o coeficiente de \textit{resplast} quando \textit{propresp} foi adicionada à regressão? Isso faz sentido?
    \item Adicione \textit{mailsyear}, o número de correspondências enviadas por ano, ao modelo. Quão grande é seu efeito estimado? Por que essa pode não ser uma boa estimativa do efeito causal dos envios sobre as respostas?
\end{enumerate}

\end{questions}
\newpage
\bibliographystyle{apalike}
\bibliography{refs}
\end{document}
